\documentclass[10pt,letterpaper]{article}
\usepackage[utf8]{inputenc}
\usepackage[T1]{fontenc}
\usepackage[hmargin=1.65cm,vmargin=1.15cm]{geometry}
\usepackage{mathptmx}

\usepackage{enumitem}
\usepackage{titlesec}
\usepackage[hidelinks]{hyperref}
\input{glyphtounicode}
\pdfgentounicode=1

\renewcommand{\familydefault}{\rmdefault}
\pagestyle{empty}
\setlength{\parindent}{0pt}
\setlength{\parskip}{1pt}
\hyphenpenalty=10000
\setlength{\emergencystretch}{2em}
\titleformat{\section}{\normalsize\bfseries\centering}{}{0pt}{}[\titlerule]
\titlespacing*{\section}{0pt}{7pt}{4pt}
\setlist[itemize]{leftmargin=14pt,itemsep=1pt,topsep=2pt,parsep=0pt}
\hypersetup{pdftitle={Juan Luis Soto Martinez | Business Central e integraciones},pdfauthor={Juan Luis Soto Martinez},pdfsubject={Curriculum vitae}}
\begin{document}
\begingroup\centering
{\fontsize{18}{21}\selectfont\bfseries Juan Luis Soto Martinez}\par
{\normalsize Especialista en Integraciones | Dynamics 365 Business Central}\par
Monterrey, Nuevo León, México \enspace\textbar\enspace \href{mailto:jlusomart@gmail.com}{jlusomart@gmail.com}\par
\href{https://linkedin.com/in/grandeotos}{linkedin.com/in/grandeotos} \enspace\textbar\enspace
\href{https://github.com/grandeotos}{github.com/grandeotos}\par\endgroup

\section{Perfil profesional}
Especialista en integraciones con más de 4 años en desarrollo de software y backend. Autor de más de 20 extensiones en AL y responsable de conectar ERP, comercio electrónico y PIM mediante APIs y Azure. Principal punto de escalación técnica del área de sistemas y referente para incidencias complejas de Business Central. Global Tenant Administrator de Microsoft 365 / Entra ID.

\section{Experiencia profesional}
\textbf{ARTEXA México | Especialista en Integraciones}\par
Monterrey, Nuevo León \hfill Jul 2026 -- Actualidad\par
\begin{itemize}
\item Encabezo la escalación técnica de mayor nivel del área de sistemas y la integración y mejora continua de sistemas empresariales, con Business Central como eje.
\item Diseño (en curso) la conciliación automática de pagos: referencia por documento, lectura de archivos bancarios y exposición de anticipos mediante API.
\item Apoyo la coordinación y la extracción de datos en la migración multiempresa de Business Central y de sistemas ERP heredados a la nube, en conjunto con una consultora externa.
\item Administro la infraestructura de Azure y el tenant de Microsoft 365 / Entra ID, superviso técnicamente el software, doy mentoría a un desarrollador junior y elaboro guías operativas para ventas y cumplimiento de pedidos.
\end{itemize}
\textbf{ARTEXA México | Desarrollador de Software Senior} \hfill Remoto · Jul 2024 -- Jun 2026
\begin{itemize}
\item Desarrollé y mantuve más de 20 extensiones de Business Central en AL (reportes, integraciones, automatización de flujos y personalización de procesos) y administré el ERP.
\item Lideré la integración backend de \textbf{ArtexaPRO}: conexión de Adobe Commerce con Business Central para pedidos, inventario y datos de productos, incluida la optimización (V2) del servicio de consulta de pedidos.
\item Diseñé una \textbf{capa de integración central (API Gateway)} que unifica Business Central, la tienda en línea y el catálogo de productos en un solo punto de acceso, simplificando las conexiones entre sistemas.
\item Detecté y corregí la causa raíz de corrupción de datos en el motor de posteo de Business Central (reservation entries e integridad fiscal), con herramientas en AL para su saneamiento.
\item Apoyé la migración de suscripciones de Azure (PAYG a CSP, en curso), la migración del correo transaccional a Azure Communication Services y la evolución de un asistente generativo en producción (catálogo, precios, logística y facturación).
\end{itemize}
\textbf{ARTEXA México | Desarrollador de Software Junior} \hfill Remoto · Jun 2023 -- Jul 2024\par
Desarrollo de extensiones en AL y personalización de reportes y facturas de Business Central; apoyo en configuración de usuarios y soporte de TI.
\par\textbf{ARTEXA México | Practicante de Desarrollo de Software} \hfill Remoto · Dic 2022 -- Jun 2023\par

\vspace{3pt}
\textbf{Clik Graphics | Desarrollador Backend} \hfill Jul 2022 -- Oct 2022\par
Desarrollo de APIs REST con NestJS, TypeScript y TypeORM; apoyo en migración de una base de datos heredada.

\vspace{3pt}
\textbf{OMRI Tech | Desarrollador Backend} \hfill Ago 2021 -- Feb 2022\par
Desarrollo y mantenimiento de una plataforma LMS y sitios web con PHP, Laravel y WordPress.

\section{Competencias técnicas}
\textbf{ERP:} Dynamics 365 Business Central, AL, codeunits, páginas, reportes, event subscribers e integraciones.\par
\textbf{Backend e integración:} Python, FastAPI, Pydantic, APIs REST, OpenAPI, Adobe Commerce y Power Automate.\par
\textbf{Cloud y datos:} Azure Container Apps, Functions, Logic Apps, DevOps, Communication Services, Entra ID; PostgreSQL, MongoDB, MySQL. \textbf{Herramientas:} Git, Docker, Linux y Postman.

\section{Formación e idiomas}
\textbf{Ingeniería en Tecnologías Computacionales} \hfill 2019 -- 2024\par
Tecnológico de Monterrey, Campus Cuernavaca.\par
\textbf{Idiomas:} Español nativo; inglés B2+/C1 (LinguaSkill).\par
\textbf{Certificación:} SCRUM Fundamentals Certified, SCRUMstudy (2021).
\end{document}
